% DELTA23 /root/five_level_six; proof draft for independent checking.

\subsection*{D23.F3. An original fifth-fibre Tate character on the gentle group}
Let $F=\widehat F_2=\langle x,y\rangle$ with $z=(xy)^{-1}$,
and let $\mathcal H=\ker\lambda\subset\widehat{GT}_0$ act by
$\alpha_h(x)=x$, $\alpha_h(y)=f_h^{-1}yf_h$.
Consider the cyclic quotient
\[
 F\twoheadrightarrow C_5\times C_2,
 \qquad x\mapsto(1,0),\quad y\mapsto(0,1),
\]
with kernel $K_{10}$, and let $N$ be the closed normal subgroup
of $F$ generated by $x^5,y^2,(xy)^{10}$.
Then $N\subset K_{10}$ and
\[
 \Pi_5=K_{10}/N
\]
is the completed fundamental group of the compact genus-two
curve $C_5:y^2=x^5-1$.  Every $h\in\mathcal H$ induces an
automorphism of $\Pi_5$ and a continuous character
\[
 \Omega_5:\mathcal H\longrightarrow
 \{u\in(\widehat{\mathbb Z}[\zeta_5])^\times:
                         u\overline u=1\}.
 \tag{D23.F3.1}
\]
On $G_{\mathbb Q^{\rm ab}}$ this is the full Tate representation
of $J_5=\operatorname{Jac}(C_5)$, in the fixed covering marking.
If $f_h\in F''$, then $\Omega_5(h)=1$.

\begin{proof}
The algebraic map
\[
 C_5\longrightarrow\mathbb P^1_t,
 \qquad t=x^5=y^2+1
\]
is Galois of degree ten over $\mathbb Q(\zeta_5)$, with deck
group $C_5\times C_2$.  Its ramification orders over $0,1,\infty$
are $5,2,10$.  Removing these fibres gives the covering of
$\mathbb P^1\setminus\{0,1,\infty\}$ specified by $K_{10}$;
filling them gives exactly the normal closure of the displayed
peripheral powers.  This identifies $\Pi_5$.  Riemann--Hurwitz
gives
\[
 2g-2=10\left(-2+(1-1/5)+(1-1/2)+(1-1/10)\right)=2,
\]
so $g=2$.

The gentle action is peripheral: each of $x,y,z$ is sent to a
conjugate of itself at multiplier one.  It is the identity on
the abelian quotient $C_5\times C_2$.  Hence it preserves
$K_{10}$ and $N$, with the same true statement for its inverse;
the resulting map on $\Pi_5$ is an automorphism.  It commutes
with the order-five deck action on $\Pi_5^{\rm ab}$, since it
fixes $x$ literally.

The induced map on the four-dimensional homology is symplectic.
Here is a direct boundary argument that specifies its multiplier.
The open cover has eight punctures: two over zero, five over one,
and one over infinity.  The two over zero are fixed because $x$
is fixed.  The five over one are fixed individually because the
conjugator $f_h$ lies in $F'\subset K_{10}$, so it has trivial
deck coordinate.  The puncture over infinity is unique.
All eight peripheral loops are therefore preserved up to conjugacy,
with exponent one.  Fill the first seven punctures but leave
the one over infinity.  The resulting completed free group has
rank four; in a geometric free basis its remaining boundary is
\[
 [a_1,b_1][a_2,b_2].
\]
Our induced automorphism preserves its conjugacy class.  In
the quotient of the lower central series of degrees two and
three, conjugation disappears and this boundary gives
$a_1\wedge b_1+a_2\wedge b_2$.  Preservation of that element
is precisely the symplectic equation.  This proof works after
every finite reduction, hence over $\widehat{\mathbb Z}$.

The order-five rotation on $C_5$ has eigenvalues
$\zeta_5,\zeta_5^2$ on $dx/y,x\,dx/y$, and the conjugate
eigenvalues on the antiholomorphic part.  Its characteristic
polynomial on integral first homology is $\Phi_5$.
Consequently this lattice is a torsion-free rank-one module over
$\mathbb Z[\zeta_5]$, a Dedekind ring.  Such a module is an
invertible ideal and its endomorphism ring is exactly
$\mathbb Z[\zeta_5]$; after completion the same statement
holds prime by prime.  The symplectic adjoint of rotation is
its inverse, so the adjoint on that endomorphism ring is complex
conjugation.  The commuting symplectic automorphisms are exactly
the norm-one units in (D23.F3.1).  Restriction and abelianization
respect composition, establishing the asserted character and
its continuity.

For the arithmetic identification use the tangential point
at $t=0$ directed toward $1$, lifted to $x=0,y=i$ with leading
coefficient one in the local parameter $x=t^{1/5}$.
All its constant choices lie in $\mathbb Q^{\rm ab}$.
The standard Galois inclusion in $\widehat{GT}_0$ is the action
on this marked base fundamental group.  Functoriality of the
specified algebraic cover, followed by filling its punctures,
therefore identifies its restricted homology action with the
Tate action of $J_5$.  The complex/etale comparison used here is
the Riemann existence theorem; for an exact modern locator see
Milne, \emph{Lectures on Etale Cohomology}, v2.21, Theorem 3.4.
No claim about a five-point braid frame enters this identification.

Finally, $F'\subset K_{10}$ implies $F''\subset K_{10}'$.
If $f_h\in F''$, the action on $F/K_{10}'$ fixes both $x$ and
$y$ and hence is the identity.  It is therefore the identity on
$K_{10}^{\rm ab}$ and on its quotient $\Pi_5^{\rm ab}$.
\end{proof}

\subsection*{D23.F4. A scalar cyclotomic-unit character for the normalization}
Put $K_5^{\rm cyc}=\ker(F\to C_5)$, with $x\mapsto1$,
$y\mapsto0$.  Its standard Schreier generators are
\[
 X=x^5,\qquad Y_j=x^j y x^{-j}\quad(0\leq j<5).
\]
For $h\in\mathcal H$ define
\[
 \kappa_5(h)=\operatorname{coeff}_{Y_1}
       (f_h\bmod (K_5^{\rm cyc})')\pmod2.
 \tag{D23.F4.1}
\]
This is a continuous homomorphism to $\mathbb F_2$.
With the cyclic covering $t=x^5$ and the usual radial path
from zero to one, its restriction to $G_{\mathbb Q^{\rm ab}}$
is the Kummer character of $\zeta_5-1$.

\begin{proof}
Every $f_h$ belongs to $F'\subset K_5^{\rm cyc}$.  The action
fixes $X$ and conjugates each $Y_j$ by an element of
$K_5^{\rm cyc}$.  It is therefore the identity on the
abelianization of that group.  The composition rule
$f_{hk}=f_h\alpha_h(f_k)$ now makes every integral coefficient
in that abelianization additive; reduction modulo two proves
the first assertion.

For the arithmetic coefficient, pull back the double cover of
the $x$-line given by
\[
 v^2=1-x/\zeta_5.
\]
Its sign character on the cyclic covering fundamental group
is one on $Y_1$, zero on the other $Y_j$, and zero on $X$.
The radial path from $x=0$ to $x=1$ lifts from $v=1$ to a
chosen square root of $1-\zeta_5^{-1}$.  The usual Galois
path cocycle is the lift of $f_h$, up to inversion according
to the convention; inversion has no effect on its sign.
Thus (D23.F4.1) is the quadratic Kummer character of
$1-\zeta_5^{-1}$.  Its ratio with $\zeta_5-1$ is $\zeta_5^{-1}$,
which has a square root in $\mathbb Q^{\rm ab}$.  The two
quadratic characters on $G_{\mathbb Q^{\rm ab}}$ coincide.
\end{proof}

\subsection*{D23.F5. A direct common fifth-fibre bit is impossible}
Let $k=\mathbb Q^{\rm ab}$ and retain the actual torsion fields
$L_3,L_4,L_5$ of DELTA\_ON\_GT AF.4--AF.6.  No continuous
quadratic character of $\operatorname{Gal}(L_5/k)$ pulls back
on $G_k$ to the common arithmetic bit
\[
 a_2|_{G_k}=a_3|_{G_k}
        =\kappa_{\sqrt[4]{3}}.
 \tag{D23.F5.1}
\]
Thus a proposed proof identifying both bits with one quadratic
character of the genuine $J_5$ Tate representation is false
already on the Galois subgroup.

\begin{proof}
The field of the nonzero character on the right is
$k(\sqrt[4]3)=L_3\cap L_4$, by E.65 and AF.6.
If it factored through $\operatorname{Gal}(L_5/k)$, Galois
correspondence would give $k(\sqrt[4]3)\subset L_5$.
But AF.6 states $L_3\cap L_4\cap L_5=k$, a contradiction.
This applies to the full Tate image, so passing to a finite
level cannot repair the contradiction.
\end{proof}

\paragraph{Scope.}
D23.F3 and D23.F4 construct actual homomorphisms on the gentle
group, with arithmetic interpretations.  D23.F1--D23.F2
construct a different, frame-induced character on full GT.
An equality between the two constructions is not part of
these statements.  The explicit quadratic twist above shows
which scalar must be included in a comparison at the chosen
normalization; it does not make such a comparison on all GT
automatic.
